\section{Keeping the same company through every calculation}
\label{sec:northstar-portable}

The first inventory purchase shows what a reproducible company calculation
must retain. Before the purchase, Northstar has six units costing 30 and owes
suppliers 20. It buys fourteen more units for 70 on credit. The journal debits
inventory and credits trade payables by 70; the resulting inventory is twenty
units costing 100, payables are 90, and cash remains 60. A balanced journal
accompanied by twenty-one units of inventory is wrong under the stipulated
unit cost of 5. The quantity change and the monetary posting must therefore be
checked against the same purchase and opening state.

Delivery follows purchase. Its four journal lines recognise the receivable and
revenue together with cost of sales and the reduction in inventory. For the
first half-year, delivery of fifteen units leaves five units costing 25.
Separating the revenue and inventory effects into independently committable
events could expose a state in which the customer has the goods but the seller
still reports them. Treating delivery as one economic event prevents that
intermediate inconsistency. Collections then precede the payments they fund.
This order establishes a feasible aggregate history; it does not recover the
dates of individual invoices.

The portable company example retains twenty-two successive states and journal
steps: ten economic events and a closing entry in each half-year. Each step
identifies its prior state, economic inputs, accounting template and resulting
state. The corrected June statements and December statements refer to this same
chain. The December closing state, with cash 74.6 and equity 90.6, is also the
explicit opening of the 2025 projection. Consequently, changing a historical
purchase or tax payment cannot leave the future opening silently unchanged.
These are proposed histories; publishing them as candidates does not post them
to an accounting book.

The rate comparison checks a different connection. Recall that $x$ is the
floating share of the unchanged principal of 100. For an imposed floating rate
$r$, the fixed part earns interest at 5\%, so annual interest is
\begin{equation}
 I(x;r)=100(1-x)0.05+100xr=5+(100r-5)x.
 \label{eq:northstar-portable-rate}
\end{equation}
At 7\%, increasing $x$ from zero to one adds 2 to interest; at 9\%, it adds 4.
Under the stipulated 20\% tax rule, profit and cash therefore fall by 1.6 or
3.2 respectively. The second rate supplies a useful control: two authoring
frontends might agree because they share an erroneous constant. Independent
arithmetic must recover the changed slope from the changed rate. The reference
exchanges retain the exact expression for every $x\in[0,1]$, including its
dependence across cash, profit, assets and equity. They attach no probabilities
to the endpoints or to intermediate debt allocations.

The opening also determines what kind of forecast this is. The corrected
December and interim information becomes public in March 2025, after the
January origin of the scenario. Using it is an archival conditional exercise:
given that corrected opening, what follows from these imposed decisions and
rates? It cannot establish what an investor could have forecast in January.
An investor-information request at that earlier cutoff rejects the later
disclosure dependency. Separately, a management policy cannot read a future
realised collection merely because the simulation scheduler already stores it.
An imposed scenario assumption and a later observed outcome have different
information meanings.

Partial information must survive the same chain. If the fraction of current
sales collected is missing, delivery can still determine revenue and cost of
sales. Cash, receivables and the ability to fund the proposed settlements remain
unresolved. The partial report identifies its affected lines and retains the
full debt family. If the tax method is unavailable, profit after tax, retained
earnings and the final closing entry are unresolved as well. A zero in a
temporary account would incorrectly suggest that closing had been completed.
The portable values therefore carry an explicit reason and its outstanding
obligation wherever the calculation is unavailable.

There is a similar distinction in reading a note. The separate authored phrase
``substantially all customers pay promptly'' does not imply a 90\% threshold.
An analyst may impose an all-payers assumption for a conditional calculation,
but its decision must preserve the unquantified source reading. Likewise,
interpreting a displayed 100 as nearest-unit rounding requires an explicit
rounding assumption. With ties rounded upward, that assumption permits values
in $[99.5,100.5)$; absent the assumption, the display alone supplies no such
interval. Neither decision grants permission to post a journal.

Finally, a correct purchase must survive an uncertain acknowledgement. The
request identifies Northstar, its book and branch, and the particular purchase;
it also fixes the opening state, expected revision and all input versions.
Suppose the client sends the request but receives no response. If the store
committed nothing, an exact retry may commit the purchase. If the store already
committed it, the same retry returns the original receipt without another
purchase. The client cannot distinguish these possibilities from a timeout.

Receipt lookup therefore precedes comparison with today's head, after the
original request has passed its input and authorization checks. A later,
separately identified purchase may have advanced the book, but it does not
erase the first receipt. Replaying the first request still returns that receipt.
Changing its request identifier or budget changes its full fingerprint and
conflicts with the already recorded business event. A new event presented
against an obsolete opening is instead rejected as stale. Every rejection
leaves the committed states, journals, quantity changes and receipt index
unchanged.

Review decisions, saved search progress and imports have their own owners and
receipts. Correcting the reported revenue invalidates dependent current uses,
including a selection that depended on the old set of eligible readings; it
does not rewrite committed purchases. Saving one searched region of
$[0,1]$ preserves the other region and the next position to visit. Resuming
with a larger budget updates that saved work, while a changed model requires
a different calculation. These distinctions make the complete company example
reconstructible from its inputs and record references. The operation-specific
fields and exact documentary exchanges are described in
Section~\ref{sec:corporate-exchange-reference}; they still require implementation
and actual fault testing before any claim of database durability.
